\subsection{Realization over three-manifolds with finite fundamental group}

The finite cyclic hypothesis can be removed from the surjectivity
assertion in Theorem~\ref{thm:natural-three-concordance-obstruction}.
The relevant surgery input is the relative obstruction formula of
Hambleton, Milgram, Taylor and Williams; the construction of the
homeomorphism is Galvin's relative realization construction.

\begin{theorem}\label{thm:finite-three-concordance-realization}
Let $B$ be a closed, connected, oriented real PL three-manifold
with finite fundamental group.  Then
\[
 c_B:\pi_0P_{\mathrm{TOP}}(B)\longrightarrow
                         H^2(B;\mathbb Z/2)
\]
is surjective.  Every class on the right has a representative
$F:B\times I\to B\times I$ fixing both ends pointwise.
Here $P_{\mathrm{TOP}}(B)$ has only the bottom face fixed,
and its components are ordinary topological path components.
The images of the jointly PL component map and of the definable
component map for every fixed o-minimal expansion of
$\mathbb R$ are contained in $\ker c_B$.
\end{theorem}

We use the formal relative PL difference class
$\operatorname{cs}_{\mathrm{PL}}(f)\in
H^3(X,\partial X;\mathbb Z/2)$ defined in
\eqref{eq:relative-three-base-cs}.  It records the prescribed
stable PL reductions on the entire boundary of the
five-dimensional mapping cylinder.

\begin{lemma}\label{lem:finite-sylow-relative-cs-realization}
Let $X$ be a compact, connected, oriented smooth four-manifold,
possibly with boundary.  Suppose that $\Gamma=\pi_1(X)$ is finite
and that a Sylow two-subgroup $\rho$ satisfies
\[
                    SK_1(\mathbb Z[\rho])=0.
\]
Then each $\eta\in H^3(X,\partial X;\mathbb Z/2)$ is
$\operatorname{cs}_{\mathrm{PL}}(f)$ for a homeomorphism
$f:X\to X$ fixing the entire boundary pointwise.
\end{lemma}

\begin{proof}
Set $M=X\times I$ and retain its full boundary, including
$\partial X\times I$.  Use fixed product collars, or round
the corners compatibly.  The relative normal-invariant
coordinates of
\cite[Lemmas~4.8--4.9, p.~2526]{Galvin2026CassonSullivan}
are
\begin{equation}\label{eq:finite-relative-normal-coordinates}
 [M/\partial M,G/\mathrm{TOP}]
 \cong H^2(M,\partial M;\mathbb Z/2)
                \oplus H^4(M,\partial M;\mathbb Z).
\end{equation}
The stable PL obstruction of $(0,y)$ is
$\operatorname{red}_2(y)$.  The coefficient exact sequence and
orientation give
\[
 H^4(M,\partial M;\mathbb Z)
 \xrightarrow{\operatorname{red}_2}
 H^4(M,\partial M;\mathbb Z/2)\longrightarrow
 H^5(M,\partial M;\mathbb Z)
 \xrightarrow{\,2\,}H^5(M,\partial M;\mathbb Z).
\]
Since $H^5(M,\partial M;\mathbb Z)\cong\mathbb Z$,
reduction is onto.  Choose $y$ reducing to the interval
suspension $\sigma\eta$.

We verify the vanishing of the surgery obstruction of this
particular normal invariant.  In the notation of
\cite[Theorem~A(b) and Definition~(0.6), p.~351]
{HambletonMilgramTaylorWilliams1988}, the permitted decoration
$U\subset\operatorname{Wh}(\Gamma)$ must contain
\[
 Y=\operatorname{im}\bigl(SK_1(\mathbb Z[\rho])
                       \longrightarrow\operatorname{Wh}(\Gamma)\bigr).
\]
Our hypothesis permits $U=0$, the simple decoration.
Apply that theorem to the oriented Poincar\'e pair
$(M,\partial M)$, taking the identity as the reference normal
map.  For a normal invariant $\nu$ it gives
\begin{equation}\label{eq:finite-five-dimensional-surgery-formula}
 \theta^s(\nu)=\kappa_3^s\bigl(c_*
                         \operatorname{ARF}_3(\nu)\bigr),
 \qquad c:M\longrightarrow B\Gamma.
\end{equation}
The class $\operatorname{ARF}_3$ is the degree-three homology
component of
\[
       \bigl(V_M^2\smile\nu^*k\bigr)\cap[M,\partial M],
       \qquad k=\{k_{4i+2}\}_{i\geq0}.
\]
Since $\dim M=5$, this component uses only $k_2$ and the
degree-zero Wu factor.  Its pullback vanishes for a normal
invariant with zero degree-two coordinate in
\eqref{eq:finite-relative-normal-coordinates}.  Therefore
\begin{equation}\label{eq:finite-zero-arf-realization}
                  \theta^s(0,y)=0.
\end{equation}
In particular, the integral degree-four coordinate does not
contribute to this obstruction.  This argument requires no
vanishing assertion for the whole group $L_5^s(\mathbb Z[\Gamma])$.

The simple relative TOP surgery sequence
\cite[Theorem~11.3A, p.~200]{FreedmanQuinn1990}
now realizes $(0,y)$ by a simple homotopy equivalence of triads
\[
 \phi:(N;\partial_+N,\partial_0N,\partial_-N)
   \longrightarrow
 (X\times I;X\times\{0\},\partial X\times I,X\times\{1\})
\]
whose restriction to the full boundary is the prescribed
homeomorphism.  The two end inclusions are simple homotopy
equivalences: this follows by comparing them under $\phi$
with the product inclusions, using additivity of Whitehead
torsion.  Consequently $N$ is a simple cobordism with its
labeled side product.

Finite groups are good
\cite[p.~99]{FreedmanQuinn1990}.  The compact relative
five-dimensional $s$-cobordism theorem
\cite[Theorem~7.1A, p.~101]{FreedmanQuinn1990}, in the marked
form stated in
\cite[Proposition~4.10, pp.~2526--2527]{Galvin2026CassonSullivan},
gives a product homeomorphism
\[
            p:N\xrightarrow{\cong}\partial_+N\times I
\]
which fixes the bottom and the entire labeled vertical
boundary.  Write $p_-$ for its restriction to the other end
and $\phi_\pm$ for the end markings, with the interval
coordinate suppressed.  Then
\[
                 f=\phi_+p_-\phi_-^{-1}:X\longrightarrow X
\]
All three boundary restrictions agree with the fixed side marking,
so $f|_{\partial X}=\mathrm{id}$.  To track the specified class,
put $\chi=(\phi_+\times\mathrm{id}_I)p:N\to X\times I$.
The maps $\phi$ and $\chi$ agree on the lower end and on the
entire vertical boundary.  Consequently
\[
 H=\operatorname{pr}_X\phi\chi^{-1}:X\times I\longrightarrow X
\]
is a homotopy relative to $\partial X$ from $\mathrm{id}_X$
to $f^{-1}$.  In particular, $f$ acts trivially on
$H^3(X,\partial X;\mathbb Z/2)$.

The two markings $\phi$ and $\chi$ induce the same map on
$H_1(N;\mathbb Z/2)$, because they agree on an end whose
inclusion is a homotopy equivalence.  Both carry the relative
mod-two fundamental class to that of $X\times I$.
Naturality of Poincar\'e--Lefschetz duality therefore identifies
their pullbacks on $H^4(X\times I,\partial(X\times I);\mathbb Z/2)$.
Thus the normal-invariant marking and the product marking assign
the same class $\operatorname{red}_2(y)$ to the relative
obstruction.

In the $\chi$ coordinates the lower reduction is the standard
one, the upper reduction is its pullback by $f^{-1}$, and the
side reduction is the prescribed product reduction.  The marked
mapping-cylinder identity
\cite[Lemma~4.12, p.~2527]{Galvin2026CassonSullivan} gives
\[
 \operatorname{cs}_{\mathrm{PL}}(f^{-1})
       =\sigma^{-1}\operatorname{red}_2(y)=\eta.
\]
The inverse rule for formal difference classes and the trivial
cohomology action give
\[
 \operatorname{cs}_{\mathrm{PL}}(f)
   =f^*\operatorname{cs}_{\mathrm{PL}}(f^{-1})=\eta,
\]
as required.
\end{proof}

\begin{proof}[Proof of Theorem~\ref{thm:finite-three-concordance-realization}]
We first verify the Sylow hypothesis of the lemma.  A finite
triangulation of $B$ lifts to a finite free cellular
$\mathbb Z[\Gamma]$-complex $C_*$ on its universal cover,
where $\Gamma=\pi_1(B)$.  The cover $\widetilde B$ is compact,
simply connected and oriented.  Simple connectivity and
three-dimensional Poincar\'e duality give
\[
 H_0(\widetilde B;\mathbb Z)
   =H_3(\widetilde B;\mathbb Z)=\mathbb Z,
 \qquad
 H_1(\widetilde B;\mathbb Z)
   =H_2(\widetilde B;\mathbb Z)=0.
\]
The deck group acts trivially on $H_3$, since the orientation
is pulled back from $B$.  Hence there is an exact sequence
of $\mathbb Z[\Gamma]$-modules
\begin{equation}\label{eq:finite-three-periodic-extension}
 0\longrightarrow\mathbb Z\xrightarrow{\iota}C_3
 \longrightarrow C_2\longrightarrow C_1
 \longrightarrow C_0\xrightarrow{\epsilon}\mathbb Z
 \longrightarrow0,
\end{equation}
with trivial modules at both ends.  Splice copies of this
extension, joining consecutive blocks by
$\iota\epsilon:C_0\to C_3$.  Exactness at a join follows
from
\[
 \ker(\iota\epsilon)=\ker\epsilon=\operatorname{im}d_1,
 \qquad
 \operatorname{im}(\iota\epsilon)
     =\operatorname{im}\iota=\ker d_3.
\]
This constructs a four-periodic free resolution and proves
that $\Gamma$ has periodic cohomology.

For a finite group with periodic cohomology, its Sylow
two-subgroups have trivial integral $SK_1$; this is the
specific arithmetic conclusion recorded in the last
introductory paragraph of
\cite[p.~352]{HambletonMilgramTaylorWilliams1988}.
The corresponding computations for cyclic and generalized
quaternion two-groups are also given in
\cite[Introduction, Example~2, p.~14; Example~14.4,
p.~333]{Oliver1988Whitehead}.
We use this assertion for the Sylow subgroup alone.

Choose a smoothing of $B$ compatible with its PL structure
\cite[Essay~I, p.~39]{KirbySiebenmann1977}
and set $X=B\times I$ with the product smoothing.
Given $a\in H^2(B;\mathbb Z/2)$, its interval suspension
lies in $H^3(X,\partial X;\mathbb Z/2)$.
Lemma~\ref{lem:finite-sylow-relative-cs-realization}
realizes this class by a homeomorphism $F:X\to X$
fixing both ends.  The normalization of $c_B$ in
Theorem~\ref{thm:natural-three-concordance-obstruction}
then gives $c_B([F])=a$.
The same theorem gives both kernel containments, with
ordinary top-free components and the stipulated joint PL
and real definable parameter conventions.
\end{proof}

The obstruction homomorphism is therefore onto. Determining
its kernel and the two comparison images remains a further question.
